\documentclass{article}
\usepackage{graphicx} % Required for inserting images
\usepackage[english]{babel}

\usepackage[letterpaper,top=2cm,bottom=2cm,left=3cm,right=3cm,marginparwidth=1.75cm]{geometry}

\usepackage{amsmath}
\usepackage{graphicx}
\usepackage[colorlinks=true, allcolors=blue]{hyperref}

\title{Data 501 Literature Review I}
\author{Ashton Frese, Karen Reyes González, Wilbur Elbouni}
\date{February 10, 2026}

\begin{document}

\maketitle 

\title{Paper Review: \textit{The Effects of Macroeconomic Shocks on Sector-Specific Returns}}

\section{Paper Overview}

\noindent \textbf{Citation:}  
B. T. Ewing, S. M. Forbes, and J. E. Payne, ``The effects of macroeconomic shocks on sector-specific returns,'' \textit{Applied Economics}, vol. 35, pp. 201--207, 2003.

\vspace{0.5em}

\noindent \textbf{Problem Addressed:}  
The paper investigates how unexpected changes (``shocks'') in key macroeconomic variables affect the returns of major sector-specific stock market indices in the United States. Instead of focusing solely on an aggregate market index, the authors examine heterogeneous responses across sectors.

\vspace{0.5em}

\noindent \textbf{Motivation:}  
With the growing use of sector and market indexes as performance benchmarks and as the basis for index funds, options, derivatives, and other structured products, there is a practical need to understand:

\begin{itemize}
    \item which macroeconomic variables drive sector-level returns,
    \item how large these effects are,
    \item and how long they persist following a macroeconomic shock.
\end{itemize}

\noindent The authors state that investors, fund managers, and policymakers benefit from a clearer view of how sector-specific indices respond to macroeconomic news, especially when sector returns may respond differently to a market shock than an individual stock or aggregate market measure.

\section{Methodology Summary}

\noindent \textbf{Data:}  
The study uses monthly U.S. data from 1988:1 to 1997:7, focusing on five Standard \& Poor's sector stock price indices:

\begin{itemize}
    \item Capital Goods (LNCAP),
    \item Financials (LNFI),
    \item Industrials (LNIN),
    \item Transportation (LNTR),
    \item Utilities (LNUT).
\end{itemize}

\noindent These sector indices are analyzed alongside three macroeconomic variables:

\begin{itemize}
    \item \textbf{Real Output:} The index of coincident indicators (LNCOINC), used as a proxy for real output.
    \item \textbf{Monetary Policy:} The effective Federal Funds Rate (FYFF), capturing the stance of monetary policy.
    \item \textbf{Market Risk:} The spread between a Baa-rated corporate bond yield and the 10-year Treasury yield (RISK), representing a market-wide risk premium.
\end{itemize}

\noindent All price/index series (except interest rates) are transformed to natural logs and then first-differenced to obtain growth rates/returns. Interest rate spreads (RISK) are used at levels.

\vspace{0.5em}

\noindent \textbf{Techniques and Models:}  
The authors estimate a series of six-variable Vector Autoregressions (VARs), each including the five sector returns and one macroeconomic variable:

\begin{itemize}
    \item VAR1: sector returns + output growth (DLNCOINC),
    \item VAR2: sector returns + change in Fed Funds Rate (DFYFF),
    \item VAR3: sector returns + risk premium (RISK).
\end{itemize}

\noindent Lag length is selected primarily using information criteria (AIC, SBC) and likelihood ratio tests, a single lag is chosen for the reported models. Stationarity is checked with standard unit root tests, and no cointegration is detected among the variables, motivating a VAR in differences for most series.

%This part seems pretty complex, kinda hard to give the full context in a simpler manner tho, bunch of previous methods are mentioned, maybe take that part out and only focus on GIRFs ??
%We should also omit a bunch of this and put it in five-year-old terms, lmao
\noindent A key methodological choice in the paper is the use of 
\textbf{generalized impulse response functions} (GIRFs), following Koop et al. (1996) 
and Pesaran and Shin (1998). GIRFs avoid the ordering assumptions required by 
traditional Cholesky-based orthogonalized impulse responses. They:

\begin{itemize}
    \item do not rely on an arbitrary variable ordering,
    \item allow contemporaneous correlation among shocks,
    \item and produce responses that are invariant to reordering.
\end{itemize}


\noindent Generalized impulse responses are then computed for a one-standard-deviation shock to each macroeconomic variable, tracing the response of each sector index over a five-month horizon.

\section{Key Contributions}

\noindent \textbf{Main Findings:}

\begin{itemize}
    \item \textbf{Output Shocks (Real Activity):}  
    A positive shock to real output growth (coincident index) generates an immediate \textit{positive} and statistically significant impact on all five sector indices. However, this is followed by a rapid reversal and overshooting, with returns briefly turning negative before the effect dissipates. The magnitude of the initial impact and the extent of overshooting differ by sector, with capital goods and financials showing stronger initial effects, and utilities exhibiting relatively low initial sensitivity but more pronounced overshooting.
    
    \item \textbf{Monetary Policy Shocks (Fed Funds):}  
    An unanticipated tightening of monetary policy (increase in Fed Funds Rate) causes an immediate \textit{negative} and statistically significant response in all sector returns. The largest initial impacts are observed for financials and capital goods, reflecting their interest-rate sensitivity and leverage. The effects for most sectors decay within roughly two months, with utilities showing a small overshooting pattern.
    
    \item \textbf{Risk Premium Shocks:}  
    A positive shock to the risk premium (widening Baa–Treasury spread) leads to a \textit{positive} and significant initial response for all sectors, with the strongest effect in utilities, followed by financials. The dynamics again differ by sector, where utilities and transportation show distinctive patterns, consistent with their regulatory constraints and exposure to input costs (e.g., energy prices).
\end{itemize}

\vspace{0.5em}

\noindent \textbf{Innovations and Insights:}

\begin{itemize}
    \item The paper explicitly challenges what it maintains has been an implicit presumption in the literature, namely, ``that all sectors respond identically to macro shocks"
    \item With their use of generalized impulse response analysis as opposed to order-dependent orthogonalized impulse responses, the study demonstrates more robust findings regarding the overall effect of macro shocks.
    \item The study finds that output, monetary policy, and the risk premium matter for sector-specific returns, each affecting different sectors in unique ways.
\end{itemize}

\section{Critical Reflection}

\noindent \textbf{Strengths:}

\begin{itemize}
    \item \textbf{Sector-Level Focus:}  
    The decision to work with sector indices instead of the aggregate index allows directly for the modeling of the question of which sectors of the market are more or less sensitive to macroeconomic news.
    
    \item \textbf{Clear Macroeconomic Narrative:}  
    The paper's approach to isolate real activity, monetary policy, and risk premium provides an easily interpretable story for how those factors are propagated into the returns within the sectors.
    
    \item \textbf{Robust Time-Series Framework:}  
    Using VARs and using generalized responses is appropriate from a methodological point of view; this is especially true because of the contemporaneous correlation between the sectors.
\end{itemize}

\vspace{0.5em}

\noindent \textbf{Limitations and Assumptions:}

\begin{itemize}
    \item \textbf{Data Window and Frequency:}  
    Shocks are statistically determined (variance innovations of the VAR-residuals) rather than through an explicit event database (for instance, specific FOMC announcements, geopolitical events) crises, or labelled economic shocks). This limits the ability to map responses to concrete real world events in a dashboard or narrative.
    
    \item \textbf{Implicit Event Treatment:}  
    Shocks are defined statistically (innovations in the VAR residuals) rather than via an explicit event database (e.g., identified FOMC announcements, geopolitical crises, or labelled economic shocks). This limits the ability to map responses to concrete real-world events in a dashboard or narrative.
    
    \item \textbf{No Forecasting Evaluation:}  
    The research explores dynamic response and qualitative behaviour rather than formal forecasting of performance. There is no comparison of predictive models (e.g., ARIMA, GARCH, LSTM) or evaluation metric such as
    
    \item \textbf{No Visual Analytics Component:}  
    The results are communicated using static plot displays do not explore interactive or visual analytics techniques, which are central to the focus of our project’s dashboard.
\end{itemize}

\section{Connection to Our Project}

\noindent This paper is highly relevant in multiple ways to our capstone project: framing the problem, selection of variables, and methodological design. It also highlights important gaps that our work aims to address.

\vspace{0.5em}

\noindent \textbf{Conceptual Alignment:}  
Our current project aims at the development of a 'one-stop' interactive dashboard to perform analysis and forecasting of sector performance in the U.S. equity market with respect to macroeconomic and geopolitical events. Ewing et al. illustrate that:

\begin{itemize}
    \item that macroeconomic shocks to real activity, monetary policy, and risk premium have measurable and heterogeneous impacts across sectors,
    \item and that sector-level analysis conveys sensitivities which might have been obscured in an aggregate index. 
\end{itemize}

\noindent This directly supports our decision to work with sector ETFs in comparison to a broad market index alone, such as XLK, XLE, XLF, XLV, XLI. It also motivates our emphasis on cross-sector comparison in the dashboard.

\vspace{0.5em}

\noindent \textbf{Influence on Variable and Event Design:}  
The use of the Federal Funds Rate as a core monetary policy variable informs our inclusion of the Effective Federal Funds Rate and our research question on easing vs.\ tightening cycles. Ewing et al.\ treat monetary policy shocks as changes in the Fed Funds Rate and show that sectors such as Financials and Capital Goods are especially sensitive. We extend this by:

\begin{itemize}
    \item examining how sector ETF returns and volatility co-move with the Federal Funds Rate over a longer horizon (2005–2026),
    \item and contrasting behaviour during distinct economic policy regimes (e.g., zero lower bound, 2015–2018 normalization, 2022–2023 tightening).
\end{itemize}

\noindent Similarly, their treatment of real activity and risk premium as distinct channels motivates our interest in augmenting the Fed Funds data with additional macro indicators in future extensions of the dashboard.

\vspace{0.5em}

\noindent \textbf{Methodological Inspiration and Differences:}  
Methodologically, the paper motivates a shock-oriented, dynamic view of sector behaviour. While we are not planning on implementing VARs or generalized impulse responses directly, we adopt several of their underlying ideas:

\begin{itemize}
    \item \textbf{Event Windows as an Analogue to Impulse Responses:}  
    A similar study to this is conducted by Ewing et al., where they analyze the path of responses in different sectors to an unanticipated macroeconomic innovation. Similarly, while working on our assigned project, we attempt to approximate a similar function by specifying event windows, say 30 days before and after a macroeconomic or geopolitical event as per information provided in Wikidata, and then computing their respective returns and volatilities.

    \item \textbf{From Descriptive Dynamics to Forecasting:}  
    Unlike the focus of Ewing et al. on the descriptive dynamics of the reactions of sectors to shocks, we consider the predictive performance of these models as well. Our metrics of performance of these models (RMSE, MAE) for the respective hold-out period include shock regimes as well as stable regimes.

\end{itemize}


\vspace{0.5em}

\noindent \textbf{Visualization and Dashboard Design:}  
The impulse response plots in the paper reveal a consistent visual structure where we observe the evolution of different sectors in response to the same shock over time. Our approach to visualizing the framework is similar but based on the previously mentioned perspective.

\begin{itemize}
    \item \textbf{Event Overlay Timeline:} Similar to their sectoral impulse responses, our dashboard overlays major events (and policy changes) impacting the price and volatility series related to the sector ETFs at Wikidata. allowing users to visually inspect heterogeneous sector reactions.
    
    \item \textbf{Drill-Down View:} Inspired by their focus on dynamic response windows, users can click an event (e.g., a Fed policy change or a globally significant crisis) to zoom into a specified window and compare sector returns and volatility trajectories. This effectively provides an interactive equivalent to the impulse response plots.
\end{itemize}

\vspace{0.5em}

\noindent \textbf{Extension of Scope:}  
Finally, our work extends the scope of Ewing et al.\ in several ways:

\begin{itemize}
    \item \textbf{Data:} From monthly S\&P sector indices (1988–1997) to daily U.S. sector ETFs (2005–2026), covering the 2008 financial crisis, the COVID-19 pandemic, and the recent tightening cycle.
    \item \textbf{Events:} From statistically defined shocks to an explicit, curated list of macroeconomic and geopolitical events, filtered via Wikidata sitelinks for global importance.
    \item \textbf{Models:} From VAR-based descriptive dynamics to a comparative forecasting framework (ARIMA, GARCH, LSTM) evaluated under different market regimes.
    \item \textbf{Interface:} From static econometric figures to an interactive Plotly/Dash dashboard that operationalizes the idea of sector-specific responses to macro shocks for exploratory analysis.
\end{itemize}

\noindent In summary, Ewing et al.\ provide a sector-level, macro-shock perspective that directly supports the conceptual foundation of our capstone project. Their findings justify our sector focus, our emphasis on monetary policy and macro events, and our use of event windows, while also highlighting important methodological and visualization gaps that our interactive dashboard aims to address.

\end{document}